\documentclass[11pt,a4paper]{article}

% ---------------------------------------------------------------------------
% Packages
% ---------------------------------------------------------------------------
\usepackage[utf8]{inputenc}
\usepackage[T1]{fontenc}
\usepackage{lmodern}
\usepackage[margin=2cm]{geometry}
\usepackage{graphicx}
\usepackage{booktabs}
\usepackage{multirow}
\usepackage{amsmath,amssymb}
\usepackage{siunitx}
\usepackage{caption}
\usepackage{float}
\usepackage[hidelinks]{hyperref}
\usepackage{xcolor}
\usepackage{enumitem}
\usepackage{titlesec}

% ---------------------------------------------------------------------------
% Global settings
% ---------------------------------------------------------------------------
% The twelve figure PDFs (fig01_..._fig12_...) must sit in a graphs/ folder next to
% this file (or compile from the folder that contains them, e.g. the simulation
% output folder).  Nothing else is needed beyond a standard TeX Live install.
\graphicspath{{graphs/}{./graphs/}{figures/}{./figures/}}
\captionsetup{font=small,labelfont=bf,skip=6pt}
\sisetup{detect-all}
\setlist{nosep,leftmargin=1.4em}

\setlength{\parskip}{0.5em}
\setlength{\parindent}{0pt}
\titlespacing*{\section}{0pt}{1.5em}{0.8em}
\titlespacing*{\subsection}{0pt}{1.2em}{0.6em}

% ---------------------------------------------------------------------------
% One consistent figure macro.
%   * Every figure sits alone on its own page (float [p]).
%   * Every figure is scaled to the SAME bounding box:
%         width  = \linewidth
%         height <= 0.78\textheight
%     so nothing ever overflows the page.
% ---------------------------------------------------------------------------
\newcommand{\reportfig}[3]{%
  \clearpage
  \begin{figure}[p]
    \centering
    \includegraphics[width=\linewidth,height=0.78\textheight,keepaspectratio]{#1}
    \caption{#2}
    \label{#3}
  \end{figure}
  \clearpage
}

\title{\bfseries Simulation 1: Median-Based Committee Selection under Manipulation\\[0.4em]
\large A Detailed, Plain-Language Figure Report (corrected edition)}
\author{}
\date{\today}

\begin{document}
\maketitle

\begin{abstract}
This report explains, in simple language, the twelve figures produced by
Simulation~1. The simulation asks a straightforward question: \emph{if some
voters lie about their preferences, can they push a candidate they like into
the winning committee?} To answer it, we compare three ways of choosing a
committee --- Evaluative Voting, the \(k\)-Median Rule, and the
\(k\)-Median-Shapley Rule --- against two kinds of attacks (the Cutoff attack
and the Bottom attack), under six different models of what voters' honest
opinions look like, and for eight committee sizes (\(k = 2, \dots, 9\)). The
electorate is always \(100\) voters, there are \(m = 10\) candidates, and
every combination is repeated many times to get stable numbers. Every figure
below is explained from scratch: what is on each axis, what each colour and
line means, what the shape of the curve tells us, and what the overall
conclusion is.

\medskip
\noindent\textbf{Note on the figure set.} The figures in this report are the
\(m = 10\) build, with committee sizes \(k = 2, \dots, 9\) and centred
\(5\)-point smoothing on the drawn curves. If an older \(m = 5\) figure set
(\(k = 2,3,4\), no smoothing) is used, several numbers below --- especially
the median-rule tipping points and the IC ordering --- will differ. Always
match the report to the figure set that was actually rendered.

\medskip
\noindent\textbf{Where the numbers come from.} Figures and cell-level
statistics are quoted from the \(m = 5\) and \(m = 7\) runs at \(1{,}000\)
trials per cell, and from the audited \(m = 10\) verification run at \(200\)
trials per cell (the run whose audit reports \textsc{all checks passed}).
Individual \(m = 10\) curve values move by a point or two between runs, but
the ranges, the orderings and the qualitative statements below are stable.

\end{abstract}

\tableofcontents
\clearpage

% ===========================================================================
\section{What Is This Simulation About?}
% ===========================================================================

Imagine a small election. There are \(100\) voters and \(10\) candidates.
Every voter honestly scores every candidate from \(0\) to \(5\). We then pick a committee of size
\(k\) (for example, \(k = 5\)) using some rule. That is the \emph{honest
committee}.

Now suppose a group of voters decides to lie. They all agree to give the
highest possible score to one specific candidate --- their \emph{target} ---
and the lowest possible score to everybody else. They hope this pushes the
target into the committee.

The simulation asks:
\begin{itemize}
  \item How many liars does it take before the committee changes?
  \item How many liars does it take before the target actually gets in?
  \item Does the answer depend on which voting rule we use?
  \item Does the answer depend on what the honest voters' opinions look like?
  \item Does the answer depend on how big the committee is?
\end{itemize}

The twelve figures are different ways of looking at those answers.

% ===========================================================================
\section{The Building Blocks}
% ===========================================================================

% ---------------------------------------------------------------------------
\subsection{The Six Preference Models}
% ---------------------------------------------------------------------------

A ``preference model'' is just a recipe for generating honest voter opinions.
We use six of them, because real electorates can look very different.

\begin{table}[H]
  \centering
  \caption{The six ways we generate honest voter opinions.}
  \label{tab:cultures}
  \begin{tabular}{@{}p{4.2cm}p{10.5cm}@{}}
    \toprule
    Model & In plain words \\
    \midrule
    Mallows (\(\varphi = 0.7\)) &
    Voters mostly agree on a single ``correct'' ranking, but each voter makes
    a few random swaps. Dispersion is controlled. \\
    Plackett--Luce &
    Voters pick candidates one at a time from a weighted lottery. Popular
    candidates are more likely to be picked first. \\
    2D Euclidean &
    Each voter and each candidate is a dot on a square map. Voters like
    candidates who are close to them. \\
    IAC (Impartial Anonymous Culture) &
    For each candidate independently, the whole electorate draws a random
    distribution of scores. There is no shared structure. \\
    P\'olya Urn (\(\alpha = 0.2\)) &
    Voters copy each other. Each new vote is more likely to look like votes
    already cast. This creates clusters of similar opinions. \\
    Impartial Culture (IC) &
    Every voter gives every candidate a completely random score from
    \(\{0,1,2,3,4,5\}\). The most disordered model possible. \\
    \bottomrule
  \end{tabular}
\end{table}

% ---------------------------------------------------------------------------
\subsection{The Three Voting Rules}
% ---------------------------------------------------------------------------

A ``rule'' is a recipe for turning all the voters' scores into a committee.

\begin{itemize}
  \item \textbf{Evaluative Voting.}
        Add up every candidate's scores across all voters. Pick the \(k\)
        candidates with the biggest totals. This is like a popular vote ---
        the more total points you get, the better.
  \item \textbf{\(k\)-Median Rule.}
        For each candidate, find the median score (the middle score when all
        voters are lined up from lowest to highest). Pick the \(k\) candidates
        with the highest medians. This cares about the typical voter, not the
        total, so a handful of extremists cannot move it very much.
  \item \textbf{\(k\)-Median-Shapley Rule.}
        Use the median, but instead of scoring a candidate by their own
        median alone, score them by their \emph{average marginal
        contribution}: take every possible group (coalition) of candidates
        the candidate could join, look at the median of the group's combined
        score with and without them, and average the difference over all
        such groups. This is the Shapley value of the median-based game. It
        rewards candidates who lift up the whole committee, not just
        themselves.
\end{itemize}

% ---------------------------------------------------------------------------
\subsection{The Two Attacks}
% ---------------------------------------------------------------------------

The liars always use the same ballot: score~\(5\) for their target, \(0\) for
everyone else. What changes is \emph{which} candidate they target.

\begin{itemize}
  \item \textbf{Cutoff attack.} The target is the candidate ranked \(k+1\)
        by the honest electorate --- the first candidate who just missed the
        committee. This is the ``easiest'' candidate to push in, because they
        are already close to the boundary.
  \item \textbf{Bottom attack.} The target is the candidate ranked last
        (rank~\(m\)). This is the hardest candidate to push in, because they
        start from the very bottom.
\end{itemize}

% ---------------------------------------------------------------------------
\subsection{The Three Metrics}
% ---------------------------------------------------------------------------

For each trial we record three numbers.

\begin{itemize}
  \item \textbf{Change rate (\%)} --- did the committee change at all?
        \(100\%\) means it changed in every trial; \(0\%\) means it never
        changed.
  \item \textbf{Success rate (\%)} --- did the target actually enter the
        committee? This is what the liars really care about.
  \item \textbf{Average overlap} --- what fraction of the honest committee
        survived? \(1.0\) means nothing changed; \(0.0\) means the whole
        committee was replaced.
\end{itemize}

% ===========================================================================
\section{How to Read Every Figure}
% ===========================================================================

Before we look at each figure, here is a small cheat-sheet.

\begin{itemize}
  \item \textbf{Horizontal axis (\(x\))} is always ``manipulators as a
        percentage of the \(100\)-voter electorate''. It goes from \(1\%\) to
        \(49\%\). So moving right means more liars.
  \item \textbf{Vertical axis (\(y\))} is always a percentage rate
        (\(0\)--\(100\%\)) or an overlap (\(0\)--\(1\)). Higher means the
        attack is doing more damage.
  \item \textbf{Colour} tells you the rule:
        \textcolor[HTML]{D55E00}{orange} \(=\) Evaluative Voting,
        \textcolor[HTML]{0072B2}{blue} \(=\) \(k\)-Median,
        \textcolor[HTML]{009E73}{green} \(=\) \(k\)-Median-Shapley.
  \item \textbf{Line style and marker} tell you the committee size~\(k\).
        All eight sizes are drawn with the \emph{same} line width
        (\(1.7\)\,pt); only the dash pattern and the marker change.
        \(k=2\) is the finest dotted line with circle markers, \(k=3\) the
        dashed line with squares, \(k=4\) the solid line with triangles, and
        so on up to \(k=9\) (dash--dot--dot, plus markers). The legend at the
        bottom of each figure shows the exact mapping.
  \item \textbf{Panel layout.} Figures~1--6, 10, 11, 12 are \(2 \times 3\)
        grids (six preference models). Figure~7 is a heatmap (one panel per
        rule). Figure~8 is a single tall panel with the legend at lower right.
        Figure~9 is a \(2 \times 2\) grid (rows = attacks, columns = metrics).
  \item \textbf{About the ``larger \(k\) is more manipulable'' shorthand.}
        This is true for the \emph{bottom} attack but \emph{not} for the
        \emph{cutoff} attack. Under the cutoff attack small committees are the
        more exposed ones. Wherever this report says ``larger \(k\)'' or
        ``smaller \(k\)'' about manipulability, it will say which attack it
        means. See Section~\ref{sec:contradiction}.
\end{itemize}

% ===========================================================================
\section{Figure-by-Figure Walkthrough}
% ===========================================================================

% ---------------------------------------------------------------------------
\subsection{Figure 1 --- Cutoff Attack: How Often Does the Committee Change?}
% ---------------------------------------------------------------------------

\reportfig{fig01_cutoff_change_rate.pdf}%
{Cutoff attack: committee change rate across the six preference models.
Each panel shows one culture. Colour \(=\) rule, line style and marker \(=\)
committee size \(k\). The \(y\)-axis is the percentage of trials in which the
winning committee differs from the honest committee. Curves are drawn with a
centred \(5\)-point moving average (display only).}{fig:cutoff_change}

\textbf{What is plotted.} For each preference model (one panel), for each rule
(colour), for each committee size (line style), the figure shows the percentage
of trials in which the committee changed, as we increase the number of liars
from \(1\%\) to \(49\%\).

\textbf{What to look for.}
\begin{itemize}
  \item \textbf{Orange lines shoot up almost immediately.} Under Evaluative
        Voting, even a handful of liars is enough to change the committee.
        The reason is simple: adding up scores is very sensitive to extreme
        ballots.
  \item \textbf{Blue lines rise more slowly.} The \(k\)-Median Rule looks at
        the middle voter, not the sum. A single extreme ballot does not move
        the median much.
  \item \textbf{Green lines are the slowest.} The Shapley rule spreads the
        influence of any one candidate across many possible coalitions.
  \item \textbf{Under the cutoff attack, \emph{small} \(k\) is the more
        exposed one --- clearly so for evaluative voting, where the \(k=2\)
        curve is never below the \(k=9\) curve in any dataset (\(m=5\), \(7\)
        and \(10\)); at \(m=10\) it is strictly above it at \(61\%\) of
        intensities when the cultures are averaged (the other \(39\%\) are
        ties, mostly where both curves have saturated). For the median rules
        the ordering across \(k\) is much weaker and partly non-monotone: at
        \(m=10\) the \(k\)-Median \(k=9\) curve is above the \(k=2\) curve at
        \(53\%\) of intensities and below it at \(37\%\). This is the
        \emph{opposite} of what happens under the bottom attack (see
        Figure~\ref{fig:bottom_change} and
        Section~\ref{sec:contradiction}).}
  \item \textbf{Impartial Culture is the most fragile.} In a completely random
        electorate, no rule can hold its ground for long.
  \item \textbf{Mallows and Plackett--Luce are the most stable.} When voters
        mostly agree, the honest signal is strong and liars have less
        influence. Averaged over the first \(15\%\) of manipulators at
        \(m = 5\), the evaluative change rate runs \(90.5\%\) under Impartial
        Culture but only \(53.5\%\) under Mallows and \(48.9\%\) under
        Plackett--Luce.
\end{itemize}

\textbf{Take-away.} If the question is ``how easily can a small group change
the committee?'', the answer is: \emph{very easily under evaluative voting,
much harder under median rules, and hardest under Shapley.}

% ---------------------------------------------------------------------------
\subsection{Figure 2 --- Cutoff Attack: How Often Does the Target Get In?}
% ---------------------------------------------------------------------------

\reportfig{fig02_cutoff_success_rate.pdf}%
{Cutoff attack: target success rate. Same layout as
Figure~\ref{fig:cutoff_change}, but the \(y\)-axis now records the percentage
of trials in which the targeted cutoff candidate actually enters the
committee.}{fig:cutoff_success}

\textbf{What is plotted.} The same setup as Figure~\ref{fig:cutoff_change}, but
now we do not ask ``did anything change?''. We ask ``did the \emph{target}
get in?''.

\textbf{What to look for.}
\begin{itemize}
  \item \textbf{For evaluative voting, Figures~1 and~2 are identical.}
        Under evaluative voting, the only way the committee changes is if the
        target swaps in, so change rate equals success rate exactly.
  \item \textbf{For the \(k\)-Median rule, change and success are also
        identical under the cutoff attack.} The reason is structural: under
        the cutoff attack, the target is the single candidate the honest
        committee excludes, so ``the committee changed'' and ``the target got
        in'' describe the same event. Across the full \(m = 10\) bottom-vs-cutoff
        audit, this equivalence holds in \(100\%\) of cells.
  \item \textbf{For \(k\)-Median-Shapley, a small gap appears.} The Shapley
        rule can reshuffle the committee without admitting the target --- up
        to roughly \(35\) percentage points in the worst single cell of the
        \(m = 10\) dataset (mean gap \(4.7\)~pp).
\end{itemize}

\textbf{Take-away.} Under the cutoff attack, ``change'' and ``success'' are
almost the same question for the two non-Shapley rules; the Shapley rule is
the only one where churn is not a perfect proxy for admitting the target.

% ---------------------------------------------------------------------------
\subsection{Figure 3 --- Bottom Attack: How Often Does the Committee Change?}
% ---------------------------------------------------------------------------

\reportfig{fig03_bottom_change_rate.pdf}%
{Bottom attack: committee change rate. The target is now the worst-ranked
candidate. Layout and encoding follow Figures~\ref{fig:cutoff_change}
and~\ref{fig:cutoff_success}.}{fig:bottom_change}

\textbf{What is plotted.} Exactly Figure~\ref{fig:cutoff_change}, but now the
liars try to push in the \emph{worst-ranked} candidate instead of the first
runner-up.

\textbf{What to look for.}
\begin{itemize}
  \item \textbf{All curves shift right relative to the cutoff attack.}
        Pushing the worst candidate in is harder, so it takes more liars.
  \item \textbf{Under the bottom attack, \emph{large} \(k\) is the more
        exposed one.} Averaged over cultures at \(m = 10\), the \(k = 9\)
        curve lies above the \(k = 2\) curve at \(86\%\) of intensities under
        evaluative voting, \(94\%\) under \(k\)-Median and \(88\%\) under
        Shapley. The pattern weakens in the most disordered models --- in
        Impartial Culture the same share falls to \(29\%\), \(18\%\) and
        \(80\%\) respectively, because there both curves saturate early. This
        is the reverse of the cutoff attack and is the source of the apparent
        contradiction in earlier drafts
        (Section~\ref{sec:contradiction}).
  \item \textbf{The curves are monotone.} In the \(1{,}000\)-trial runs the
        drawn curves rise smoothly (at \(m=5\), \(99\%\) of consecutive steps
        are non-decreasing for \(k\)-Median, \(96\%\) for Shapley, with a
        largest single dip of \(3\)~pp). In a \(200\)-trial pilot the same
        curves show \(\pm3\)--\(7\)~pp wobbles; those are Monte Carlo noise,
        not structure.
  \item \textbf{The Shapley rule is still the most robust.} Its change rate
        stays below \(k\)-Median's in \(91\%\) of cells and below evaluative
        voting's in \(76\%\) of cells in the \(m = 10\) data --- the ordering
        holds, but it is a tendency rather than an absolute law.
\end{itemize}

\textbf{Take-away.} The bottom attack is harder than the cutoff attack, and
in this regime bigger committees are easier to flip.

% ---------------------------------------------------------------------------
\subsection{Figure 4 --- Bottom Attack: How Often Does the Target Get In?}
% ---------------------------------------------------------------------------

\reportfig{fig04_bottom_success_rate.pdf}%
{Bottom attack: target success rate. The \(y\)-axis records the percentage of
trials in which the bottom-ranked candidate enters the committee.}{fig:bottom_success}

\textbf{What is plotted.} The most important figure for the paper's main
claim.

\textbf{What to look for.}
\begin{itemize}
  \item \textbf{Orange curves (evaluative voting)} rise fast and reach
        \(100\%\) success well before \(49\%\) liars --- earliest in the most
        disordered electorate (Impartial Culture, at \(15\%\) liars) and
        latest in the best-structured ones (Mallows at \(37\%\),
        Plackett--Luce at \(43\%\)). Note the direction: the disordered models
        fall first, not last.
  \item \textbf{Blue curves (\(k\)-Median)} rise more slowly. In the \(m=10\)
        run, their \(50\%\) tipping points span roughly \(5\%\)
        (Impartial Culture, \(k=9\)) to \(42\%\) (Mallows, \(k=2\)).
  \item \textbf{Green curves (\(k\)-Median-Shapley)} are the lowest of all. In
        the \(m = 10\) build, six of the \(48\) (culture, \(k\)) cells never
        reach \(50\%\) success inside the simulated range, and all six are at
        \(k = 2\) or \(k = 3\). Larger committees eventually let the target
        through.
  \item \textbf{Culture changes the height, not the order --- with one
        caveat.} The \(k\)-Median rule beats evaluative voting in \(90.5\%\)
        of the bottom-attack cells. All \(9.5\%\) of exceptions sit at low
        manipulation (\(80\) cells at \(1\)--\(5\%\) liars, \(131\) at
        \(6\)--\(15\%\), \(12\) at \(16\)--\(30\%\), none beyond \(30\%\)),
        and never by more than \(9\)~pp.
\end{itemize}

\textbf{Take-away.} This is the headline result: \emph{median-based rules are
substantially harder for liars to exploit, especially for small committees.}

% ---------------------------------------------------------------------------
\subsection{Figure 5 --- Cutoff Attack: How Much of the Committee Survives?}
% ---------------------------------------------------------------------------

\reportfig{fig05_cutoff_overlap.pdf}%
{Cutoff attack: average overlap with the honest committee. Overlap is the
fraction of the honest committee that survives in the post-manipulation
committee, averaged over trials.}{fig:cutoff_overlap}

\textbf{What is plotted.} Overlap is a smoother metric than ``did anything
change?''. If the honest committee is \(\{A,B,C\}\) and the post-attack
committee is \(\{A,B,D\}\), the overlap is \(2/3 \approx 0.67\).

\textbf{What to look for.}
\begin{itemize}
  \item \textbf{Orange curves collapse quickly.} Under evaluative voting,
        overlap drops from \(1.0\) towards the ``changed committee''
        \emph{ceiling}
        \[
          \frac{k-1}{k} \;=\; 1 - \frac{1}{k},
        \]
        which is the largest overlap you can have when the committee is not
        the honest committee. This is \(0.500\) at \(k=2\), \(0.667\) at
        \(k=3\), and rises towards \(1\) as \(k \to m\).
  \item \textbf{The true disjoint-committee floor} (when the new committee
        shares nobody with the honest one) is
        \[
          \frac{\max(0,\,2k - m)}{k},
        \]
        which is zero for all \(k \le m/2\) and only becomes positive at
        \(k = 6, 7, 8, 9\) in the \(m = 10\) build. Observed minima in the
        \(m = 10\) cutoff-attack data are \(0.138\) (\(k = 2\), Shapley) and
        \(0.185\) (\(k = 2\), \(k\)-Median) --- even the harsher bottom attack
        bottoms out at \(0.125\) --- comfortably above the naive
        ``\(1 - k/n\)'' guess that appeared in an earlier draft.
  \item \textbf{Blue and green curves stay near \(1.0\) for a long time.}
        Under Mallows, the median rules keep a mean overlap above \(0.9\)
        until about \(23\%\) liars, and \(93\%\) of the individual cells
        within the first \(20\%\) of manipulators still have overlap
        \(\ge 0.9\).
  \item \textbf{Impartial Culture is the harshest environment.} Its
        culture-averaged curve for the median rules falls from \(0.73\) at
        \(10\%\) liars to \(0.57\) at \(49\%\), the lowest of the six models,
        and individual (culture, \(k\)) cells drop below \(0.5\) early ---
        \(6.7\%\) of IC cells do so within \(1\)--\(15\%\) liars, the minimum
        being \(0.282\). (The averaged curve itself never falls below
        \(0.5\).)
\end{itemize}

\textbf{Take-away.} The median rules protect the whole committee structure,
not just the target slot.

% ---------------------------------------------------------------------------
\subsection{Figure 6 --- Bottom Attack: Change Rate vs.\ Success Rate Together}
% ---------------------------------------------------------------------------

\reportfig{fig06_bottom_dual_metric.pdf}%
{Bottom attack: change rate (solid) vs.\ target success rate (dashed). The
shaded band spans the range across committee sizes \(k = 2,\dots,9\); the
solid/dashed lines are the means.}{fig:bottom_dual}

\textbf{What is plotted.} Both metrics on the same axes. The vertical gap
between a solid line and its dashed counterpart is ``wasted churn'' ---
committee changes that did not admit the target.

\textbf{What to look for.}
\begin{itemize}
  \item \textbf{For evaluative voting, solid and dashed lines sit on top of
        each other.} No wasted churn.
  \item \textbf{For the median rules, the two lines separate widely.} The
        gap is largest at \emph{small} \(k\) and shrinks as \(k\) grows. At
        \(k = m - 1 = 9\) the gap is exactly zero, because at that size the
        bottom-attack target is the single candidate the honest committee
        excludes --- so ``the committee changed'' and ``the target got in''
        are literally the same event.
  \item \textbf{Typical gap sizes.} Averaged over the \(m = 10\) data, the
        change-minus-success gap under the bottom attack is
        \(23.4\)~pp at \(k = 2\) for \(k\)-Median and \(31.4\)~pp at \(k = 2\)
        for Shapley, falling to \(0.0\)~pp at \(k = 9\) for both. The worst
        single cells reach \(67\)~pp (\(k\)-Median) and \(73\)~pp (Shapley);
        one clean example is P\'olya Urn, \(k = 2\), Shapley, at \(43\%\)
        manipulation: \(85\%\) change versus \(15\%\) success, a \(70\)~pp
        gap.
  \item \textbf{The gap is attack-dependent.} Under the \emph{cutoff} attack,
        the \(k\)-Median gap is exactly zero in \(100\%\) of cells --- the
        median rule's churn under the cutoff attack is always a win for the
        target.
\end{itemize}

\textbf{Take-away.} Robustness in this simulation is churn that does not
admit the target, and it is concentrated at small \(k\) under the bottom
attack.

% ---------------------------------------------------------------------------
\subsection{Figure 7 --- Bottom-Target Success Map}
% ---------------------------------------------------------------------------

\reportfig{fig07_success_heatmap.pdf}%
{Bottom-target success map. Rows are culture \(\times\) committee size,
columns are manipulation bins; one panel per rule. Colour encodes the
bottom-target success rate (\%). Dark red \(= 100\%\), white \(= 0\%\).}{fig:heatmap}

\textbf{What is plotted.} A colour-coded grid. Each row is one (culture, \(k\))
pair; each column is a range of manipulation intensities.

\textbf{What to look for.}
\begin{itemize}
  \item \textbf{Left panel (evaluative voting).} The grid fills in from the
        left. The \(10\)--\(15\%\) column is \emph{mid-scale} orange-red
        (mean success \(46\%\); only \(31\%\) of its rows are above \(80\%\)),
        not deep red. Most rows reach \(80\%\) or more from the
        \(30\)--\(35\%\) column onward (\(96\%\) of them), and every row is at
        \(100\%\) from \(35\)--\(40\%\) onward.
  \item \textbf{Centre panel (\(k\)-Median).} A slower gradient; small \(k\)
        rows stay pale much longer.
  \item \textbf{Right panel (\(k\)-Median-Shapley).} The palest of all ---
        \(25\) of its \(48\) rows are still below \(80\%\) success in the
        last column, \(13\) of them below \(50\%\), and all of those sit at
        small \(k\) (\(k \le 8\), mostly \(k \le 5\)).
  \item \textbf{Reading the rows.} Each culture is a block of eight rows,
        ordered \(k=2\) (top of block) to \(k=9\) (bottom of block).
        \emph{Under the bottom attack}, the larger-\(k\) rows are the darker
        ones --- the reverse of the cutoff attack.
\end{itemize}

\textbf{Take-away.} The heatmap summarises the whole simulation at a glance:
saturation moves left-to-right in the order evaluative \(\to\) \(k\)-Median
\(\to\) Shapley.

% ---------------------------------------------------------------------------
\subsection{Figure 8 --- Tipping Points}
% ---------------------------------------------------------------------------

\reportfig{fig08_tipping_points.pdf}%
{Manipulation level (\% of the electorate) at which the bottom target first
reaches a \(50\%\) success rate. Further right \(=\) more robust. Open
markers on the right edge indicate the threshold was never reached inside
the simulated \(1\)--\(49\%\) range.}{fig:tipping}

\textbf{What is plotted.} A single tall panel with one row per
(culture, \(k\)) pair. On each row there are three markers --- one per rule
--- placed at the manipulation level where success first reaches \(50\%\).
Markers pushed to the right edge and drawn with a hollow centre mean ``never
reached''.

\textbf{What to look for.}
\begin{itemize}
  \item \textbf{Orange circles cluster to the left.} In the \(m = 10\) build,
        tipping points span roughly \(2\)--\(32\%\).
  \item \textbf{Blue squares sit in the middle.} Tipping points of roughly
        \(5\)--\(42\%\).
  \item \textbf{Green triangles sit furthest right.} Tipping points of
        roughly \(7\)--\(49\%\), and six of the \(48\) rows are hollow ---
        all at \(k = 2\) or \(k = 3\).
  \item \textbf{The horizontal spread between the three markers is the
        robustness gain.} On almost every row, the order from left to right
        is orange, blue, green.
\end{itemize}

\textbf{Take-away.} This is the cleanest single-picture summary of the rule
ordering.

% ---------------------------------------------------------------------------
\subsection{Figure 9 --- Aggregate Summary}
% ---------------------------------------------------------------------------

\reportfig{fig09_aggregate_summary.pdf}%
{Aggregate summary: mean over the six preference models. Rows are the two
attacks; columns are the two metrics. Line style and marker encode committee
size \(k\).}{fig:aggregate}

\textbf{What is plotted.} A \(2 \times 2\) grid (rows \(=\) attacks, columns
\(=\) metrics) with the culture differences averaged away. This strips out
the noise and shows the underlying shape.

\textbf{What to look for.}
\begin{itemize}
  \item \textbf{Top row} \(=\) cutoff attack. \textbf{Bottom row} \(=\)
        bottom attack. The bottom-attack curves are always shifted right ---
        the bottom attack is harder.
  \item \textbf{Left column} \(=\) change rate. \textbf{Right column} \(=\)
        success rate. The right column curves are lower for the median rules.
  \item \textbf{The bottom-right panel} is the money shot: evaluative voting
        averages \(92\%\) success at \(30\%\) liars, \(99\%\) at \(35\%\) and
        \(100\%\) at \(49\%\), while the Shapley rule is still at \(81\%\)
        overall at \(49\%\) --- \(52\%\) at \(k = 2\), \(99\%\) at \(k = 9\).
  \item \textbf{Fan shape within each colour} \(=\) effect of committee size.
        Under the bottom attack larger \(k\) sits higher; under the cutoff
        attack the fan flips direction.
\end{itemize}

\textbf{Take-away.} Averaged over cultures, the median advantage is still
large and consistent.

% ---------------------------------------------------------------------------
\subsection{Figure 10 --- Change vs.\ Success Gap}
% ---------------------------------------------------------------------------

\reportfig{fig10_change_vs_success_gap.pdf}%
{Bottom attack: committee change rate minus target success rate (percentage
points). Zero \(=\) the target enters whenever the committee changes;
positive \(=\) committee changed but the target failed.}{fig:gap}

\textbf{What is plotted.} The vertical gap from Figure~\ref{fig:bottom_dual}
on its own, as a function of manipulation.

\textbf{What to look for.}
\begin{itemize}
  \item \textbf{Orange line sits on zero.} For evaluative voting, every change
        helps the liars.
  \item \textbf{Blue and green lines form humps.} In the \(m = 10\) data the
        per-curve humps peak at a median of \(47\)~pp (\(k\)-Median) and
        \(37\)~pp (Shapley), most of them in the \(40\)--\(60\)~pp band
        (\(43\%\) and \(36\%\) of curves). The worst single cells reach
        \(67\)~pp (\(k\)-Median) and \(73\)~pp (Shapley).
  \item \textbf{The humps are largest in the disordered models, not
        smallest.} This is the reverse of an earlier draft's claim. Mallows
        gives by far the smallest hump in all three datasets (mean peak
        \(2.5\)--\(8.4\)~pp, against \(14\)--\(48\)~pp for the other five
        models); Impartial Culture has the largest mean peak at \(m=5\)
        (\(22.3\)~pp) and \(m=7\) (\(34.4\)~pp), and at \(m=10\) the largest
        peaks are shared between Impartial Culture, P\'olya Urn and 2D
        Euclidean (\(48\)--\(50\)~pp). Intuitively, a structureless electorate
        is the easiest one for the median rules to reshuffle without helping
        the target.
  \item \textbf{Humps shrink as \(k\) grows.} In the \(m = 10\) data the
        largest gaps sit at \(k = 2\)--\(3\) (\(67\)~pp for \(k\)-Median,
        \(73\)~pp for Shapley) and the gap is exactly zero at
        \(k = m - 1 = 9\).
\end{itemize}

\textbf{Take-away.} The figure explains \emph{why} the median rules are more
robust: they can absorb the attack without letting it through, and this
mechanism is strongest for small committees and for the most disordered
electorates.

% ---------------------------------------------------------------------------
\subsection{Figure 11 --- Median Advantage}
% ---------------------------------------------------------------------------

\reportfig{fig11_median_advantage.pdf}%
{Bottom attack: robustness gain of the median rules over evaluative voting.
Positive \(=\) median rules keep the bottom target out more often than
evaluative voting.}{fig:advantage}

\textbf{What is plotted.} Direct rule-versus-rule comparison: how much more
often does the median rule block the target, compared to evaluative voting?
The \(y\)-axis is
\[
  \text{success}(\text{evaluative}) - \text{success}(\text{median}),
\]
in percentage points.

\textbf{What to look for.}
\begin{itemize}
  \item \textbf{The advantage peaks in the mid-range of manipulation.}
        In the \(m = 10\) data the Shapley-over-evaluative advantage peaks
        between \(14\%\) (at \(k = 9\)) and \(36\%\) (at \(k = 2\)) liars.
        This is the region where the median rules are most useful.
  \item \textbf{The advantage declines after the peak} because evaluative
        voting saturates at \(100\%\) success while the median rules
        eventually catch up.
  \item \textbf{Small \(k\) gives the biggest advantage.} Solid dark lines
        (small \(k\)) sit highest.
  \item \textbf{Shapley is better than \(k\)-Median.} In each colour, the
        dashed line (Shapley) sits above the solid line (\(k\)-Median).
  \item \textbf{Impartial Culture gives the \emph{smallest} advantage here.}
        This is the opposite direction from Figure~\ref{fig:gap}. Both are
        correct: IC has the largest \emph{absolute} churn gap, but the
        smallest \emph{rule-to-rule} advantage, because even evaluative
        voting is already saturated at \(100\%\) by the time the median rules
        fall short.
\end{itemize}

\textbf{Take-away.} Median rules win the most exactly where it matters ---
in the mid-range of manipulation, at small committee sizes.

% ---------------------------------------------------------------------------
\subsection{Figure 12 --- Crossing vs.\ Committee Size}
% ---------------------------------------------------------------------------

\reportfig{fig12_crossing_vs_committee_size.pdf}%
{How much manipulation is needed to elect the bottom candidate? \(x\)-axis:
committee size \(k\). \(y\)-axis: manipulation level at which the bottom
target reaches \(50\%\) success. Higher \(=\) harder to manipulate, so
higher is better.}{fig:crossing}

\textbf{What is plotted.} A compact six-panel summary, one panel per culture.
On each panel, the three curves show the \(50\%\) tipping point (from
Figure~\ref{fig:tipping}) as a function of committee size.

\textbf{What to look for.}
\begin{itemize}
  \item \textbf{All curves slope downward.} Under the \emph{bottom} attack,
        larger committees are easier to manipulate. (In the \(m = 10\) data
        every one of the \(108\) tipping-point curves is non-increasing in
        \(k\).)
  \item \textbf{Orange is uniformly the lowest curve.} Evaluative voting is
        the easiest to manipulate at every \(k\).
  \item \textbf{Green is uniformly the highest.} Shapley is the hardest to
        manipulate at every \(k\).
  \item \textbf{The vertical gap shrinks as \(k\) grows.} In the \(m = 10\)
        data the Shapley-minus-evaluative gap falls from \(\sim24\)~pp at
        \(k = 2\) to \(\sim16\)~pp at \(k = 9\).
  \item \textbf{Impartial Culture is the flattest panel for the two
        non-Shapley rules.} There the evaluative and \(k\)-Median curves move
        by less than one percentage point per step (\(0.9\)~pp and \(0.7\)~pp
        at \(m = 10\), against \(2.5\)--\(5.5\)~pp in the other panels),
        because in a random electorate the bottom target is admitted early
        for every \(k\). Only the Shapley curve still spans a wide range
        there (\(34\%\) at \(k = 2\) down to \(7\%\) at \(k = 9\)).
\end{itemize}

\textbf{Take-away.} Committee size is a trade-off: small committees are
harder to manipulate under the bottom attack, and median rules give the most
protection when the committee is small.

% ===========================================================================
\section{The Committee-Size Contradiction, Resolved}
\label{sec:contradiction}
% ===========================================================================

Earlier drafts contained two statements that looked contradictory:

\begin{itemize}
  \item ``small \(k\) is easier to manipulate'' (from the cutoff curves), and
  \item ``larger \(k\) is easier to manipulate'' (from the bottom curves).
\end{itemize}

Both are correct, but only when each is scoped to a particular attack.

\begin{itemize}
  \item \textbf{Cutoff attack.} The target is the candidate ranked \(k+1\).
        For small \(k\), that candidate is close to the middle of the
        ranking and needs only a little push to enter; for large \(k\), the
        target is deeper in the honest ranking and harder to move. So under
        the cutoff attack, \emph{small committees are the more exposed ones}.
  \item \textbf{Bottom attack.} The target is the candidate ranked \(m\).
        For small \(k\), the committee is small and the target is far away;
        for large \(k\), the committee is large and the target only has to
        overtake \(m - k\) competitors. So under the bottom attack,
        \emph{large committees are the more exposed ones}.
\end{itemize}

Every ``larger \(k\)'' or ``smaller \(k\)'' statement in this report is
therefore made with the attack named explicitly.

% ===========================================================================
\section{The Story the Figures Tell Together}
% ===========================================================================

Putting all twelve figures side by side, the following messages come out.

\begin{enumerate}
  \item \textbf{The rule ordering is almost always preserved.}
        In the \(m = 10\) audit, the ordering
        \[
          k\text{-Median-Shapley}
          \;>\;
          k\text{-Median}
          \;>\;
          \text{Evaluative Voting}
        \]
        holds in roughly \(90\)--\(100\%\) of success cells (worst pairings
        break in \(4\)--\(10\%\) of cells, and the breaks sit at low
        manipulation intensities, below about \(15\%\) liars). It is a strong
        tendency, not an absolute law.

  \item \textbf{Evaluative change equals success exactly.} In \(100\%\) of
        cells, both attacks, all three datasets, the change rate and the
        success rate of evaluative voting are identical (gap \(0.0\)~pp).

  \item \textbf{The change--success gap depends on the attack and on \(k\).}
        Under the bottom attack the gap is large at small \(k\) (up to
        \(\sim70\)~pp) and \emph{exactly zero at \(k = m-1\)}, because at
        that size the target is the only candidate excluded by the honest
        committee. Under the cutoff attack the gap is exactly zero for
        \(k\)-Median in \(100\%\) of cells.

  \item \textbf{The bottom attack is harder than the cutoff attack.} The
        bottom-attack curves shift right relative to cutoff for the success
        metric in \(100\%\) of cells (every culture, every rule); for the
        change metric the shift holds in \(98\)--\(100\%\) of cells.

  \item \textbf{Culture matters.} Impartial Culture is the harshest
        environment for every rule; Mallows is the most forgiving, followed
        by Plackett--Luce. Averaged over the first \(15\%\) of manipulators at
        \(m = 5\), the evaluative change rate is \(90.5\%\) under Impartial
        Culture against \(53.5\%\) under Mallows and \(48.9\%\) under
        Plackett--Luce.

  \item \textbf{The Shapley rule is the strongest.} Across the twelve
        figures, the \(k\)-Median-Shapley rule provides the most protection,
        especially at small \(k\) and intermediate manipulation intensities.
\end{enumerate}

% ===========================================================================
\section{Figure Index}
% ===========================================================================

\begin{table}[H]
  \centering
  \caption{Index of all figures produced by Simulation~1. Each figure is
  placed alone on its own page, at the same width and with the same caption
  style, so the reader can flip through them quickly.}
  \label{tab:figindex}
  \begin{tabular}{@{}clp{6.5cm}@{}}
    \toprule
    \# & File & Short description \\
    \midrule
    1  & \texttt{fig01\_cutoff\_change\_rate.pdf}          & Cutoff attack, change rate \\
    2  & \texttt{fig02\_cutoff\_success\_rate.pdf}         & Cutoff attack, success rate \\
    3  & \texttt{fig03\_bottom\_change\_rate.pdf}          & Bottom attack, change rate \\
    4  & \texttt{fig04\_bottom\_success\_rate.pdf}         & Bottom attack, success rate \\
    5  & \texttt{fig05\_cutoff\_overlap.pdf}               & Cutoff attack, average overlap \\
    6  & \texttt{fig06\_bottom\_dual\_metric.pdf}          & Bottom attack, change vs.\ success \\
    7  & \texttt{fig07\_success\_heatmap.pdf}              & Bottom-target success heatmap \\
    8  & \texttt{fig08\_tipping\_points.pdf}               & Tipping points \\
    9  & \texttt{fig09\_aggregate\_summary.pdf}            & Mean over cultures \\
    10 & \texttt{fig10\_change\_vs\_success\_gap.pdf}      & Change minus success gap \\
    11 & \texttt{fig11\_median\_advantage.pdf}             & Median advantage over evaluative \\
    12 & \texttt{fig12\_crossing\_vs\_committee\_size.pdf} & Tipping point vs.\ committee size \\
    \bottomrule
  \end{tabular}
\end{table}

\end{document}
